\section{Synthetic evaluation: sample error and physical structure}
\label{sec:synthetic}
\subsection{A known nonlinear field and distinct fitting experiments}
Let $I_*=\SI{100}{\ampere}$, $\psi_*=\SI{0.1}{\weber}$, and
$L_*=\psi_*/I_*=\SI{1}{\milli\henry}$. With $x=i_d/I_*$, $y=i_q/I_*$,
and $\coenergy=\psi_*I_*\normalizedcoenergy$, use
\begin{equation}
 \normalizedcoenergy=x+0.3x^2+0.5y^2
      -s(0.01x^4+0.02y^4)-0.03x^2y^2+0.02xy^2.
 \label{eq:synthetic-energy}
\end{equation}
Here $s=1$ is the baseline and $s=3$ strengthens d/q saturation. PM flux,
nonlinear self-saturation, and reciprocal cross-saturation are all present.
The normalized flux components are
\begin{align}
 \psid/\psi_*&=1+0.6x-0.04s x^3-0.06xy^2+0.02y^2,\\
 \psiq/\psi_*&=y-0.08s y^3-0.06x^2y+0.04xy,
\end{align}
with Hessian
\begin{equation}
 \Ldiff/L_*=
 \begin{bmatrix}
 0.6-0.12s x^2-0.06y^2&-0.12xy+0.04y\\
 -0.12xy+0.04y&1-0.24s y^2-0.06x^2+0.04x
 \end{bmatrix}.
\end{equation}
On $[-1,1]^2$, even at $s=3$, both diagonals are at least $0.18$ and
the off-diagonal magnitude at most $0.16$. Strict diagonal dominance proves
positive definiteness with eigenvalues at least $0.02$. This local polynomial
should not be extrapolated to arbitrarily large currents.

The experiment treats reconstructed flux samples as data with independent
normalized Gaussian noise of standard deviation $0.005$ per component. It
does not inject resistance or angle errors; those are the companion's
experiments. Draw scattered points uniformly on $[-1,1]^2$ using seed
20261006. Use 200 points for the noisy baseline, 50 for sparse sampling,
200 with five percent outliers for the outlier case, and 200 with $s=3$ for
stronger saturation. Outliers add Gaussian contamination of standard deviation
$0.1$ to both components of selected samples. No assertion about real sensor
noise is implied by this convenient synthetic distribution.

\subsection{Compare structural classes with the same spline family}
All candidates use degree-four B-splines with clamped endpoints at $-1,1$
and simple interior knots $-0.5,0,0.5$. Independent fits use two full tensor
bases. Parity fits use an even d basis and an odd q basis. The common potential
uses eight d basis functions and four even q basis functions: 32 coefficients,
31 after gauge fixing. Noiseless exact field recovery on a $13\times13$ sample
grid is checked separately; the chosen polynomial is representable in this
basis. This favorable approximation setting isolates noise and structure,
and limits conclusions about arbitrary machine maps.

Choose $\lambda$ among $0,10^{-6},10^{-5},10^{-4},10^{-3}$ using a fixed
80/20 split and validation flux error, then refit all samples with the selected
weight. Independent flux fits penalize their second derivatives, while the
potential fit penalizes third derivatives of co-energy, as described earlier.
These penalties target comparable flux smoothness but are not identical
optimization problems. The split is shared across candidates within each
case. One additional potential fit uses eight iteratively reweighted
least-squares steps with scalar Huber threshold $0.01$ for outliers.

Evaluate on a $25\times25$ grid over $[-0.95,0.95]^2$, distinct from the
training samples. Report componentwise flux RMSE in $\psi_*$ units, Hessian
entry RMSE in $L_*$ units, parity residual RMS in $\psi_*$ units, and curl
RMS in $L_*$ units. The evaluation avoids the exact outer boundary; it cannot
establish boundary or extrapolation reliability. Potential derivatives are
analytic basis derivatives, not numerical differences of interpolated flux.
\begin{table}[tbp]
\centering\scriptsize\begin{tabular}{llrrrr}
\toprule
Case & Fit & Flux RMSE & Hessian RMSE & Parity RMS & Curl RMS\\
\midrule
Noisy & Independent & 2.43e-03 & 1.40e-02 & 1.5e-03 & 2.2e-02 \\
Noisy & Parity & 1.81e-03 & 1.11e-02 & 0.0e+00 & 1.6e-02 \\
Noisy & Potential & 1.50e-03 & 1.24e-02 & 0.0e+00 & 0.0e+00 \\
Sparse & Independent & 6.80e-03 & 3.40e-02 & 4.7e-03 & 4.3e-02 \\
Sparse & Parity & 3.89e-03 & 1.79e-02 & 0.0e+00 & 2.5e-02 \\
Sparse & Potential & 2.87e-03 & 1.12e-02 & 0.0e+00 & 0.0e+00 \\
Outliers & Independent & 1.40e-02 & 8.26e-02 & 9.2e-03 & 1.0e-01 \\
Outliers & Parity & 1.03e-02 & 6.34e-02 & 0.0e+00 & 8.2e-02 \\
Outliers & Potential & 7.39e-03 & 4.74e-02 & 0.0e+00 & 0.0e+00 \\
Outliers & Potential robust & 1.50e-03 & 9.55e-03 & 0.0e+00 & 0.0e+00 \\
Strong saturation & Independent & 2.47e-03 & 1.98e-02 & 1.8e-03 & 1.7e-02 \\
Strong saturation & Parity & 1.85e-03 & 1.96e-02 & 0.0e+00 & 2.3e-02 \\
Strong saturation & Potential & 1.06e-03 & 7.91e-03 & 0.0e+00 & 0.0e+00 \\
\bottomrule
\end{tabular}

\caption{A single seeded illustrative experiment. Independent and parity
fits use flux B-splines; potential fits use co-energy gradients. Structural
zeros hold by construction. This is not an RBF benchmark or a statistical
ranking across many machines and sampling realizations.}
\label{tab:fit-metrics}
\end{table}
\begin{figure}[tbp]
\centering
\begin{tikzpicture}
\begin{groupplot}[group style={group size=2 by 2,horizontal sep=14mm,
 vertical sep=28mm},diagnostic axis,legend style={compact text,
 at={(.5,-.3)},anchor=north,legend columns=2}]
\nextgroupplot[title={d flux: slope along d},ylabel={$\psi_d/\psi_*$}]
 \addplot[reference quantity,strong line] table[x=y,y=dtrue]{figures/data/fit_slice.dat};
 \addlegendentry{True}
 \addplot[fitted curve] table[x=y,y=dfit]{figures/data/fit_slice.dat};
 \addlegendentry{Potential fit}
\nextgroupplot[title={q flux: slope along q},ylabel={$\psi_q/\psi_*$}]
 \addplot[reference quantity,strong line] table[x=y,y=qtrue]{figures/data/fit_slice.dat};
 \addplot[fitted curve] table[x=y,y=qfit]{figures/data/fit_slice.dat};
\nextgroupplot[title={d curvature},ylabel={$L_{dd}/L_*$}]
 \addplot[reference quantity,strong line] table[x=y,y=ddtrue]{figures/data/fit_slice.dat};
 \addplot[fitted curve] table[x=y,y=ddfit]{figures/data/fit_slice.dat};
\nextgroupplot[title={Mixed curvature},ylabel={$L_{dq}/L_*=L_{qd}/L_*$}]
 \addplot[reference quantity,strong line] table[x=y,y=dqtrue]{figures/data/fit_slice.dat};
 \addplot[fitted curve] table[x=y,y=dqfit]{figures/data/fit_slice.dat};
\end{groupplot}
\end{tikzpicture}
\caption{Baseline noisy-data potential fit at $x=-0.4$. Slopes and curvatures
come from the same fitted scalar surface. The fitted derivatives can show
visible error even when flux curves nearly coincide; parity and reciprocity
still hold structurally.}
\label{fig:fit-slice}
\end{figure}


\subsection{What the numbers do and do not establish}
In the baseline, the potential fit has normalized flux RMSE about $0.00150$,
versus $0.00243$ for independent fits and $0.00181$ for parity fits.
Yet its Hessian RMSE is about $0.0124$, slightly higher than the parity fit's
$0.0111$. Better pointwise prediction therefore does not automatically mean
better derivatives. Selecting regularization only on flux error can underweight
the engineering importance of differential-inductance quality.

Parity fits eliminate forbidden reflection components but retain nonzero
curl. The common potential eliminates both structurally, even in noisy and
sparse cases. This is a representation guarantee, not evidence that every
remaining prediction is true. With outliers, the ordinary potential fit
still develops a negative minimum Hessian eigenvalue of approximately
$-0.0764$ on the evaluation grid. Energy consistency and symmetry alone do
not enforce positive differential inductance. The robust potential fit
reduces flux RMSE from approximately $0.00739$ to $0.00150$ and keeps positive
sampled eigenvalues in this realization. Grid positivity still does not prove
global positivity; an explicit convexity constraint would be a separate step.

The script verifies B-spline partition of unity and analytic derivative
recurrences, gauge, gauge-fixed rank, exact noiseless flux/Hessian recovery,
and structural parity and curl of fitted candidates. It stores metrics,
selected weights, and sampled Hessian eigenvalues in
\texttt{numerics/evaluation.json}. These checks support reproducibility of
the illustration; they are not substitutes for varied-data validation.
